% coord_R2_xyz.tex
%
% Copyright (C) 2022--2026 José A. Navarro Ramón <janr.devel@gmail.com>
% Licencia del código GPLv2
% Licencia Creative Commons Recognition Non-Commercial Share-alike.
% (CC-BY-NC-SA)


\chapter{Coordenadas cartesianas ortogonales en
  \mathinhead{\symbb{R}^2}{espr2}}
Este capítulo tiene los siguientes objetivos:
\begin{enumerate}
\item Recordar el sistema de coordenadas ortogonales o rectangulares en el
  plano euclídeo $\symbb{R}^2$, como fundamento para el resto de sistemas de
  coordenadas que se estudiarán posteriormente.
\item Introducir algunas ideas de la geometría diferencial, como
  \emph{base natural}, \emph{métrica}, etc., que serán más importantes en otros
  sistemas de coordenadas.
\item Recordar la expresión del resultado de distintas operaciones de cálculo en
  este sistema.
\end{enumerate}

\section{Ideas básicas}
El vector de posición de cualquier punto del espacio $\symbb{R}^2$ en coordenadas
cartesianas es
\begin{equation}\label{eq:R2-r_rect}
  \vvv{r}
  = x\,\uvec{\i} + y\,\uvec{\j}
\end{equation}
\begin{figure}[ht]
  \pgfmathsetmacro{\EJESNEGLONG}{0.2}
  \pgfmathsetmacro{\EJESPOSLONG}{2.3}
  % Eje x
  \pgfmathsetmacro{\XNEGLONG}{\EJESNEGLONG}
  \pgfmathsetmacro{\XPOSLONG}{\EJESPOSLONG}
  % Eje y
  \pgfmathsetmacro{\YNEGLONG}{\EJESNEGLONG}
  \pgfmathsetmacro{\YPOSLONG}{\EJESPOSLONG}
  % Versores
  \pgfmathsetmacro{\VERSORMOD}{0.6}
  % Vector P
  \pgfmathsetmacro{\VECTORMOD}{2.5}
  \pgfmathsetmacro{\VECTORANG}{40}
  % Fondo
  \pgfmathsetmacro{\HORZ}{0.25}
  \pgfmathsetmacro{\VERT}{0.25}
  % 
  \centering
  \begin{tikzpicture}[%
    scale=1,
    every node/.style={black,font=\footnotesize},
    eje/.style={->},
    proyeccion/.style={black!20, line width=0.4},
    vector/.style={-{Latex[round]}, shorten >=1.2pt, line width=.8pt, black},
    versor/.style={-{Latex[round, width=4pt, length=6pt]}, line width=1pt, red!80!black},
    texto/.style={black},
    pcirculo/.style={fill=green, draw=black},
    background/.style={%
      line width=\backgrounborderwidth,
      draw=\backgroundbordercolor,
      fill=\backgroundcolor,
    },
    ]
    % -------------------------------------------------------------------------------
    % COORDENADAS
    % Origen
    \coordinate (O) at (0,0);
    % Extremo izdo del eje x e inferior del eje y
    \coordinate (xini) at (-\XNEGLONG,0);
    \coordinate (yini) at (0,-\YNEGLONG);
    \coordinate (ifin) at (\VERSORMOD, 0);
    \coordinate (jfin) at (0, \VERSORMOD);
    % Eje x
    \path[save path=\ejex] (xini) -- coordinate[pos=0.55] (xtexto) (\XPOSLONG, 0)
    coordinate (xfin);
    % Eje y
    \path[save path=\ejey] (yini) -- coordinate[pos=0.55] (ytexto) (0, \YPOSLONG)
    coordinate (yfin);
    % Versores i y j
    \path[save path=\ivector] (O) -- coordinate[pos=0.5] (itexto) (ifin);
    \path[save path=\jvector] (O) -- coordinate[pos=0.45] (jtexto) (jfin);    
    % Vector
    \path[save path=\vector] (O) -- coordinate[pos=0.5] (pmid) (\VECTORANG:\VECTORMOD)
    coordinate (P);
    % Proyección de P en los ejes
    \path[save path=\proyx] (P) |- (xfin);
    \path[save path=\proyy] (P) -| (yfin);
% -------------------------------------------------------------------------------
    % DIBUJOS
    % Proyección de P en los ejes
    \draw[proyeccion, use path=\proyx];
    \draw[proyeccion, use path=\proyy];
    % Ejes
    % x
    \draw[eje, use path=\ejex];
    \node[right] at (xfin) {$x$};
    \node[below] at (xtexto) {$x$};
    % y
    \draw[eje, use path=\ejey];
    \node[above] at (yfin) {$y$};
    \node[left] at (ytexto) {$y$};
    % Versores
    \draw[versor, use path=\ivector];
    \node[below, red!80!black] at (itexto) {\scriptsize $\uvec{\i}$};
    \draw[versor, use path=\jvector];
    \node[left, red!80!black] at (jtexto) {\scriptsize $\uvec{\j}$};
    % Vector y punto  P
    \draw[vector, use path=\vector];
    \node[above left=-3pt and -1pt] at (pmid) {$\vvv{r}$};
    \filldraw[fill=green, draw=black] (P) circle[radius=1.2pt];
    \node[above right=0pt and -5pt] at (P) {$P(x,y)$};
    
    % Origen
    \filldraw (O) circle [radius=.25pt];
%    % -------------------------------------------------------------------------------    
%    % Fondo amarillo
%    \coordinate (SW) at ($(current bounding box.south west) + (-\HORZ cm,-\VERT cm)$);
%    \coordinate (NE) at ($(current bounding box.north east) + (\HORZ cm,\VERT cm)$);    
%    \begin{scope}[on background layer]
%      \draw[background] (SW) rectangle (NE);
%    \end{scope}
  \end{tikzpicture}
  \caption{Coordenadas cartesianas rectangulares de un punto $P$ del plano
    (en verde) y vectores unitarios (en rojo). El vector de posición del punto
    es $\vvv{r} = x\,\uvec{i} + y\,\uvec{j}$.}
  \label{fig:R2-rect_ij}
\end{figure}


\section{Base natural}
En física general y en matemáticas elementales se suele preferir una base
formada por vectores unitarios para representar cualquier otro vector.
Como sabemos, en estas coordenadas se utiliza la base
$B = \set{\uvec{\i},\,\uvec{\j}}$.

Por otro lado, en geometría diferencial es importante la \emph{base natural},
que representaremos en este texto como
\begin{equation}
  B = \set{\vvv{e}_x, \vvv{e}_y}
\end{equation}
Estos vectores no tienen por qué ser unitarios, y por ello no aparecen con el
acento circunflejo, aunque como veremos, en este caso sí que lo son.

Los vectores de esta base se obtienen derivando parcialmente el vector de
posición de un punto genérico del espacio, expresado en función de las
componentes del sistema, en este caso las coordenadas $x$ e $y$, respecto de
cada una de ellas; véase la ecuación \eqref{eq:R2-r_rect} y la figura
\ref{fig:R2-rect_ij}.
Estos vectores llevan información concreta acerca de cómo varia la posición en
el espacio cuando variamos cada una de las coordenadas ---tal y como sugieren
las derivadas---, lo que será importante para obtener la \emph{métrica}.

Los vectores de la base natural y sus módulos son
\begin{subequations}
  \begin{align}\label{eq:R2-cart-ex}
    \vvv{e}_x
    &=
      \dfrac{\partial\vvv{r}}{\partial x}
      = \uvec{\i}
      = \xhat{u}_x
      \longrightarrow
    |\vvv{e}_x| = |\uvec{\i}| = 1\\
    \label{eq:R2-cart-ey}
    \vvv{e}_y
    &= \dfrac{\partial\vvv{r}}{\partial y}
      = \uvec{\j}
      = \xhat{u}_y
      \longrightarrow
    |\vvv{e}_y| = |\uvec{\j}| = 1
  \end{align}
\end{subequations}
Vemos que son vectores unitarios, aunque como se señaló antes, no tiene por qué
ser así en otros sistemas de coordenadas.

Escribimos la \emph{base natural} en cartesianas rectangulares, que coincide
con los vectores unitarios tradicionales $\uvec{\i}$ y $\uvec{\j}$
\begin{equation}\label{eq:R2-cart-base}
  B
  =\set{\vvv{e}_x, \vvv{e}_y}
  = \set{\xhat{u}_x, \xhat{u}_y}
  = \set{\uvec{\i},\, \uvec{\j}}
\end{equation}


\section{Métrica}
La métrica en cada sistema de coordenadas define la estructura del espacio, y
en particular, sin ella no se pueden medir ni distancias ni ángulos. Es como si
proporcionara un sistema de reglas y goniómetros con los cuales medir
distancias y ángulos.

De la equivalencia de la base natural con la unitaria, y del hecho de que los
vectores de la base son ortogonales entre sí, se podrá inferir que la métrica
en coordenadas cartesianas rectangulares se representará como una matriz unidad.

\subsection{Producto escalar de los vectores de la base natural}
La base natural en coordenadas cartesianas en este espacio está formada por los vectores
$\uvec{\i}$ y $\uvec{\j}$, en los ejes $x$ e $y$, respectivamente, ver figura
\ref{fig:R2-rect_ij}.
A continuación resumimos el valor de los productos escalares entre ellos, tal y como
aprendimos en cálculo vectorial fundamental en coordenadas cartesianas o en física
elemental
\begin{subequations}
  \begin{align}
    \label{eq:R2-cart-ii_prod_escalar}
    \vvv{e}_x\cdot\vvv{e}_x = \uvec{\i}\cdot\uvec{\i}
    &= |\uvec{\i}|\,|\uvec{\i}|\cos 0 = 1\\
    \label{eq:R2-cart-ij_prod_escalar}    
    \vvv{e}_x\cdot\vvv{e}_y = \uvec{\i}\cdot\uvec{\j}
    &= |\uvec{\i}|\,|\uvec{\j}|\cos\pi/2 = 0\\
    \label{eq:R2-cart-ji_prod_escalar} 
    \vvv{e}_y\cdot\vvv{e}_x = \uvec{\j}\cdot\uvec{\i}
    &= |\uvec{\j}|\,|\uvec{\i}|\cos\pi/2 = 0\\
    \label{eq:R2-cart-jj_prod_escalar}
    \vvv{e}_y\cdot\vvv{e}_y = \uvec{\j}\cdot\uvec{\j}
    &= |\uvec{\j}|\,|\uvec{\j}|\cos 0 = 1
  \end{align}
\end{subequations}

\subsection{Métrica y factores de escala}
La métrica en cualquier sistema de coordenadas se calcula utilizando los vectores de la
base natural y el producto escalar entre estos vectores. El resultado es la matriz
unidad $\mmm{I}$
\begin{equation}\label{eq:R2-cart-metrica}
  \mmm{g}_{ij}
  = \begin{pmatrix}
    g_{11} & g_{12}\\
    g_{21} & g_{22}
    \end{pmatrix}
  = \begin{pmatrix}
      \vvv{e}_x\cdot\vvv{e}_x & \vvv{e}_x\cdot\vvv{e}_y\\
      \vvv{e}_y\cdot\vvv{e}_x & \vvv{e}_y\cdot\vvv{e}_y
  \end{pmatrix}
  = \begin{pmatrix}
    1 & 0\\
    0 & 1
  \end{pmatrix}
  = \mmm{I}
\end{equation}

La raíz cuadrada de los elementos diagonales de la métrica se llaman
\emph{factores de escala} o \emph{coeficientes métricos}.
En coordenadas cartesianas rectangulares serían
\begin{align*}
  h_x &= \sqrt{g_{11}} = \sqrt{1} = 1\\
  h_y &= \sqrt{g_{22}} = \sqrt{1} = 1  
\end{align*}

Fijémonos en la métrica \eqref{eq:R2-cart-metrica}:
\begin{itemize}
\item Los elementos no diagonales son nulos. Esto se interpreta como que los
  vectores $\vvv{e}_x$ y $\vvv{e}_y$ son independientes, esto es, ortogonales,
  lo que implica que un cambio en una coordenada no afectará a la otra.
  Así, si nos desplazamos una distancia cualquiera en el eje $x$, no habremos
  recorrido ninguna distancia en el eje $y$, y viceversa, por lo que
  $h_{xy}= h_{yx}= 0$.
\item Los elementos diagonales valen la unidad. ESto implica que un cambio en
  la coordenada respectiva implica una distancia recorrida de una unidad en el
  espacio.
  En realidad, el espacio recorrido es el factor de escala correspondiente
  (la raíz cuadrada del elemento diagonal correspondiente), pero en este caso
  da lo mismo
  \begin{subequations}
    \begin{align}
      \Delta s &= h_x\, \Delta x = \sqrt{1}\, \Delta x = \Delta x\\
      \Delta s &= h_y\, \Delta y = \sqrt{1}\, \Delta y = \Delta y
    \end{align}
  \end{subequations}
  Para adquirir algo de intuición, imaginemos que nos desplazamos cualquier distancia en
  el eje $x$, digamos un metro; es fácil entender que en el espacio nos habremos
  desplazado esa misma distancia. Por tanto es lógico razonar que el factor de escala
  $h_x$ vale 1. Lo mismo podemos decir para el factor de escala en el eje $y$, $h_y=1$.
\end{itemize}

\subsection{Revisión del producto escalar}
La métrica nos permite calcular distancias. En concreto, el cuadrado del módulo de un
vector $\vvv{A}$ se mide realizando el producto escalar del vector por sí mismo.
El producto escalar se encarga de esto, y se puede ampliar mediante el cálculo
matricial, utilizando la métrica (en otros sistemas de coordenadas veremos que es más
importante que aquí, mientras que en cartesianas se puede y se suele obviar)
\[
  \vvv{v}\cdot\vvv{w}
  = \vvv{v}^\transp \mmm{g}_{ij}\, \vvv{w}
  = \vvv{v}^\transp \mmm{I}\, \vvv{w}
  = \begin{pmatrix}
    v_x &  v_y
  \end{pmatrix}
  \begin{pmatrix}
    1 & 0\\
    0 & 1
  \end{pmatrix}
  \begin{pmatrix}
    w_x\\
    w_y
  \end{pmatrix}
  =
  v_x\,w_x + v_y\, w_y
\]
Así, el cuadrado del módulo de un vector en cartesianas rectangulares es
\[
  v^2
  =
  |\vvv{v}|^2
  =
  \vvv{v}\cdot\vvv{v}
  =
  \vvv{v}^\transp \mmm{g}_{ij}\, \vvv{v}
  =
  \begin{pmatrix}
  v_x &  v_y
\end{pmatrix}
\begin{pmatrix}
  1 & 0\\
  0 & 1
\end{pmatrix}
\begin{pmatrix}
  v_x\\
  v_y
\end{pmatrix}
=
v_x^2 + v_y^2
\]
que coincide con el valor que obteníamos en el cálculo vectorial elemental.


\section{Vectores contravariantes y covariantes}
Esta sección se encuentra aquí de una manera un tanto forzada, pues los dos
tipos de vectores del título son el mismo en coordenadas rectangulares y, como
consecuencia, no es necesario hacer referencia a ellos.

Pero, adelantaremos algo de información sobre estos conceptos, para captar su
significado y entender por qué en otros sistemas de coordenadas estos vectores
son diferentes, al contrario de lo que ocurre en este sistema de coordenadas.

\subsection{Vectores contravariantes}
Los vectores \emph{viven} en un espacio vectorial $V$\, llamado
\emph{espacio tangente}.
Pensemos, por ejemplo, en un espacio curvo como puede ser el de una superficie
esférica.
No se pueden representar vectores en una superficie curva, sino en un plano.
Ese plano es en este caso el espacio tangente o espacio vectorial asociado al
punto donde \emph{viven} los vectores.

Cada punto  del espacio curvo tiene asociado un plano tangente diferente al del resto
de puntos. Nosotros sabemos realizar operaciones con vectores que
\emph{viven en el mismo espacio tangente}. Como consecuencia, en principio no podríamos
realizar operaciones entre vectores de diferentes espacios tangente\footnotemark{}.
Véase la figura \ref{fig:R2-cart_planos_tangentes}.
\footnotetext{Para poder realizar este tipo de operaciones entre dos vectores que se
  encuentren en dos espacios tangente diferentes, la geometría diferencial define el
  transporte paralelo, pero esto es algo que no nos debe preocupar por ahora.}
\begin{figure}[ht]
  % Escala
  \def\scl{1.9}
  % Origen de coordenadas
  \pgfmathsetmacro{\Ox}{0.0}
  \pgfmathsetmacro{\Oy}{0.0}
  \pgfmathsetmacro{\Oz}{0.0}
  % Centro de la esfera
  \pgfmathsetmacro{\Cx}{\Ox}
  \pgfmathsetmacro{\Cy}{\Oy}
  \pgfmathsetmacro{\Cz}{\Oz}
  % Radio de la esfera en cm
  \pgfmathsetmacro{\R}{1.0}
  % Opacidad de la esfera y los planos
  \pgfmathsetmacro{\OPACIDADESFERA}{0.9}
  \pgfmathsetmacro{\OPACIDADPLANO}{0.8}
  %
  \centering
  \begin{tikzpicture}[%
    scale=\scl,
    z= {(0cm,1cm)},
    y={(1cm,0cm)},
    x={(-0.38cm,-0.38cm)},
    line join=round,
    line cap=round,
    every node/.style={black,font=\scriptsize},
    planoP1/.style={blue!60!black, fill=blue!20, opacity=\OPACIDADPLANO},
    planoP2/.style={red!60!black, fill=red!20, opacity=\OPACIDADPLANO},
    vector/.style={-{Latex[round,width=3pt,length=4pt]}},
    ]
    % COORDENADAS
    % Origen
    \coordinate (O) at (\Ox,\Oy,\Oz);
    % Centro de la esfera
    \coordinate (C) at (\Cx,\Cy,\Cz);
    
    % 1. ESFERA
    \shade[ball color=black!20, opacity=\OPACIDADESFERA] (C) circle[radius=\R cm];
    \draw[gray!10] (C) circle[radius=\R cm];

    % 2. PLANO SUPERIOR
    \coordinate (P1) at (xyz spherical cs:radius=1,latitude=90);
    \begin{scope}[canvas is xy plane at z=1]
      % Plano en P1
      \filldraw[planoP1] (-0.5,-0.5) rectangle (0.5,0.5);
      \draw[vector] (0,0) -- (-0.3,0.3);
    \end{scope}
    % Punto P1 y texto
    \fill[black] (P1) circle (0.5pt);
    \node[below left=-2pt and -2pt] at (P1) {$P_1$};

    % 3. PLANO FRONTAL DERECHA
    \coordinate (P2) at (xyz spherical cs:radius=1,longitude=30,latitude=25);
    \begin{scope}[%
      plane origin=(P2), plane x={(0.8,1.0)}, plane y={(0.2,0)}, canvas is plane
      ]
      % Plano en P2
      \filldraw[planoP2] (-0.5,-0.5) rectangle (0.5,0.5);
      \draw[vector] (0,0) -- (0.4,-0.4);
      \draw[vector] (0,0) -- (0.25,0.4);
    \end{scope}
    
    \fill[black] (P2) circle (0.5pt);
    \node[above left=-1pt and -2pt] at (P2) {$P_2$};
  
    % Vectores en P1
    %\draw[-{Latex[round,width=3pt,length=4pt]}] (P1) -- ++(0.3, 0, 0);
    %\draw[-{Latex[round,width=3pt,length=4pt]}] (P1) -- ++(0, 0, -0.3);

%    % Vectores en P2 (ajustados para nacer exactamente de P2)
%    \draw[-{Latex[round,width=3pt,length=4pt]}] (P2) -- ++(0.05, 0.35, -0.1);
%    % \draw[-{Latex[round,width=3pt,length=4pt]}] (P2) -- ++(0.1, -0.1, -0.4);
%    \draw[-{Latex[round,width=3pt,length=4pt]}] (P2) -- ++(0.22, -0.05, -0.45);
  \end{tikzpicture}
  \caption{Espacio esférico donde se representan dos planos tangentes en dos puntos
    arbitrarios de la esfera. En el plano situado en el polo norte (azul), se representa
    un vector que \emph{vive} en ese plano. Por otro lado en el otro punto se aprecia
    un plano tangente (rosa) con dos vectores.}
  \label{fig:R2-cart_planos_tangentes}
\end{figure}

\begin{luacode*}
  -- Datos de la esfera
  esf = {
    radio = 1.0, color = "gray!10", opacidad = 1,
    sombracolor = "black!20", sombraopacidad=0.9
  }
  -- Datos del observador
  obs = {thetaD = 65, phiD = 30}
  -- 3. Matriz de puntos (tabla de tablas)
  ptos = {
    {thetaD=90, phiD=0, a=1, b=1, opacidad=0.5, color="black"},
    {thetaD=0, phiD=0, a=2, b=2, opacidad=0.5, color="red"},
    {thetaD=90, phiD=170, a=3, b=3, opacidad=0.7, color="green"},
    {thetaD=40, phiD=30, a=4, b=4, opacidad=0.7, color="blue"}
  }
  -- Planos
  planos = {}
  
  -- Añade u, v a la tabla ptos
  mEsfera.completaPuntos(esf, obs, ptos)

\end{luacode*}

{\scriptsize
  \DEBUGEsfera{esf}
  \DEBUGObservador{obs}
  \DEBUGPuntos{ptos}
  %\DEBUGPuntosAll{esf}{obs}{ptos}  
}

\begin{figure}[ht]
%  % Escala
%  \def\scl{1}
  \TIKZEsferaPlanos{esf}{obs}{ptos}
\end{figure}


%\section{Gradiente}
%Presentamos un resumen del gradiente, tal y como se presentan en un curso básico de
%cálculo.
%
%El gradiente de un campo escalar $\phi(x,y)$ es un vector
%\begin{equation}
%  \vvv{\nabla}\phi
%  = \dfrac{\partial \phi}{\partial x}\,\uvec{\i}
%  + \dfrac{\partial \phi}{\partial y}\,\uvec{\j}
%  = \begin{pmatrix}
%    \partial \phi/\partial x\\
%    \partial \phi/\partial y
%    \end{pmatrix}
%\end{equation}
%El vector gradiente de $\phi$ evaluado en un punto genérico $(x,y)$ del dominio de $\phi$
%indica la dirección en la cual el campo $\phi$ varía más rápidamente y su módulo
%representa el ritmo de variación de $\phi$ en la dirección de dicho vector gradiente.
%
%El operador gradiente se comporta como un vector bajo rotaciones, decimos esto porque
%cuando opera frente a un campo escalar, se obtiene otro formado por vectores
%\begin{equation}
%  \vvv{\nabla}
%  = \dfrac{\partial}{\partial x}\,\uvec{\i}
%  + \dfrac{\partial}{\partial y}\,\uvec{\j}
%  = \begin{pmatrix}
%    \partial/\partial x\\
%    \partial/\partial y
%    \end{pmatrix}
%\end{equation}

%\section{Divergencia}
%La divergencia de un campo vectorial $\vvv{F}=(F_x, F_y)$ es un escalar
%\begin{equation}\label{eq:R2-divergencia-cartesianas}
%  \vvv{\nabla}\cdot\vvv{F}
%  = \begin{pmatrix}
%    \partial/\partial x & \partial/\partial y      
%  \end{pmatrix}
%  \begin{pmatrix}
%    1 & 0\\
%    0 & 1
%  \end{pmatrix}
%  \begin{pmatrix}
%    F_x\\
%    F_y
%  \end{pmatrix}
%  = \dfrac{\partial F_x}{\partial x}
%  + \dfrac{\partial F_y}{\partial y}
%\end{equation}
%La divergencia mide la diferencia entre el flujo saliente y el flujo entrante de un
%campo vectorial sobre la superficie que rodea a un volumen de control, por tanto, si el
%campo tiene \emph{fuentes} la divergencia será positiva, y si tiene \emph{sumideros}
%será negativa. La divergencia mide la rapidez neta con la que se conduce el campo al
%exterior de cada punto, y en el caso de ser la divergencia idénticamente igual a cero,
%describiría el flujo incompresible del campo, llamado también campo solenoidal.
%
%\subsection{Demostración}
%Supongamos un campo vectorial $\vvv{F}$
%\[
%  \vvv{F} = F_x(x,y)\,\uvec{\i} + F_y(x,y)\,\uvec{\j}
%\]
%
%Ahora estamos interesados en investigar una propiedad del campo $\vvv{F}$ en un punto
%concreto $P(x_0,y_0)$ del espacio $\symbb{R}^2$. En particular, estamos interesados en
%estudiar si ese punto es una fuente del campo (punto en el que se crean líneas de
%campo), si es un sumidero (si desaparecen líneas), o si no ocurre ninguna de las
%dos cosas (las líneas de campo entran y salen del punto, sin más).
%
%Supongamos un pequeño cuadrado imaginario, de lado $h$ en $\symbb{R}^2$, situando el
%punto $P$ que nos interesa en el vértice inferior izquierdo, como se aprecia en la
%figura \ref{fig:R2-rect-divergencia}.
%\begin{figure}[ht]
%  % Escala
%  \def\scl{1}
%  %
%  \pgfmathsetmacro{\EJESEXTRA}{-0.4}
%  \pgfmathsetmacro{\EJESLONG}{4.0}
%  % Eje x
%  \pgfmathsetmacro{\XMLONG}{\EJESEXTRA}
%  \pgfmathsetmacro{\XPLONG}{\EJESLONG}
%  % Eje y
%  \pgfmathsetmacro{\YMLONG}{\EJESEXTRA}
%  \pgfmathsetmacro{\YPLONG}{\EJESLONG}
%  % hx, hy
%  \pgfmathsetmacro{\LONG}{2.5}
%  \pgfmathsetmacro{\HX}{\LONG}
%  \pgfmathsetmacro{\HY}{\LONG}
%  % Vector A
%  \pgfmathsetmacro{\FAMOD}{0.9}
%  \pgfmathsetmacro{\FAANG}{30}
%  % Vector B
%  \pgfmathsetmacro{\FBMOD}{0.9}
%  \pgfmathsetmacro{\FBANG}{60}
%  % Vector C
%  \pgfmathsetmacro{\FCMOD}{1.3}
%  \pgfmathsetmacro{\FCANG}{-30}
%  % Vector D
%  \pgfmathsetmacro{\FDMOD}{1.0}
%  \pgfmathsetmacro{\FDANG}{45}  
%  % Posición inferior izquierda del elemento de área
%  \pgfmathsetmacro{\PX}{0.7}
%  \pgfmathsetmacro{\PY}{0.7}
%  % Fondo
%  \pgfmathsetmacro{\HORZ}{0.25}
%  \pgfmathsetmacro{\VERT}{0.25}
%  % 
%  \centering
%  % PRIMERA IMAGEN
%  \begin{tikzpicture}[%
%    scale=\scl,
%    every node/.style={black,font=\small},
%    eje/.style={->},
%    proyeccion/.style={black!30, line width=0.3pt},
%    area/.style={fill=green!25, draw=black!20},
%    textovector/.style={black},
%    textored/.style={red!90!black},
%    textofaint/.style={black!70},
%    vector/.style={-{Latex[round,width=4pt,length=6pt]}, line width=1pt},
%    vectorcomponent/.style={-{Latex[round,width=3pt,length=5pt]}, line width=0.6pt},
%    faintvectorcomponent/.style={vectorcomponent, black!60},
%    redvectorcomponent/.style={vectorcomponent, red},
%    background/.style={%
%      line width=\bgborderwidth,
%      draw=\bgbordercolor,
%      fill=\bgcolor,
%    },
%    ]
%    % ------------------------------------------------------------------------------
%    % COORDENADAS
%    % ------------------------------------------------------------------------------    
%    % ORIGEN
%    \coordinate (O) at (0,0);
%    % EXTREMO IZQUIERDO DEL EJE X E INFERIOR DEL EJE Y
%    \coordinate (xini) at (\XMLONG cm,0);
%    \coordinate (yini) at (0,\YMLONG cm);
%    % EJES X E Y
%    \path[save path=\ejex] (xini) -- (\XPLONG cm, 0) coordinate (xfin);
%    \path[save path=\ejey] (yini) -- (0, \YPLONG cm) coordinate (yfin);
%    % PUNTO P DEL CUADRADO (punto inferior izquierdo)
%    \coordinate (P) at (\PX, \PY);
%    % PUNTOS SUPERIOR IZQUIERDA Y SUPERIOR DERECHA DEL CUADRADO
%    \path (P) -- coordinate[pos=0.5] (hxtext) ++(\HX, 0) coordinate (Pinfdcha);
%    \path (P) -- coordinate[pos=0.5] (hytext) ++(0, \HY) coordinate (Psupizda);
%    % PROYECCIÓN DE ELEMENTO DE ÁREA EN LOS EJES
%    \path[save path=\proyPx] (P) |- coordinate (Px) (xfin);
%    \path[save path=\proyPy] (P) -| coordinate (Py) (yfin);
%    \path[save path=\proyPhorx] (Pinfdcha) |- coordinate (Phorx) (xfin);
%    \path[save path=\proyPverty] (Psupizda) -| coordinate (Pverty) (yfin);
%    % TEXTO DE LAS PROYECCIONES EN LOS EJES
%    \path (Px) -- coordinate[pos=0.5] (hxtexto) (Phorx);
%    \path (Py) -- coordinate[pos=0.5] (hytexto) (Pverty);
%    % CUADRADO
%    \path[save path=\cuadrado]
%    (P) rectangle coordinate (centro) ++(\HX, \HY) coordinate (Psupdcha);
%    % PUNTOS INTERMEDIOS A (INFERIOR) Y B (SUPERIOR)
%    \path (\PX, \PY) -- coordinate[pos=0.5] (A) (Pinfdcha);
%    \path (Psupizda) -- coordinate[pos=0.5] (B) (Psupdcha);
%    % VECTORES
%    \path[save path=\FA] (A) -- ++(\FAANG:\FAMOD) coordinate (FA);
%    \path[save path=\FB] (B) -- ++(\FBANG:\FBMOD) coordinate (FB);
%    % COMPONENTES DE LOS VECTORES
%    % --- VECTOR FA
%    % RECTÁNGULO
%    \path[save path=\FArectangulo] (FA|-A) -- (FA) -- (A|-FA);
%    % COMPONENTE X
%    \path [save path=\FAx] (A) -- (FA|-A) coordinate (FAx);
%    % COMPONENTE Y
%    \path[save path=\FAy] (A) -- (A|-FA) coordinate (FAy);
%    % --- VECTOR FB
%    % RECTÁNGULO
%    \path[save path=\FBrectangulo] (FB|-B) -- (FB) -- (B|-FB);    
%    % COMPONENTE X
%    \path [save path=\FBx] (B) -- (FB|-B) coordinate (FBx);
%    % COMPONENTE Y
%    \path[save path=\FBy] (B) -- (B|-FB) coordinate (FBy);
%
%    % ------------------------------------------------------------------------------
%    % DIBUJOS
%    % ------------------------------------------------------------------------------    
%    % PROYECCIÓN DEL PUNTO P(X0,Y0) EN EL EJE X
%    \draw[proyeccion, use path=\proyPx];
%    % Punto x0
%    \fill (Px) circle[radius=0.8pt];
%    % Texto x0
%    \node[below] at (Px) {$x_0$};
%    
%    % PROYECCIÓN DEL PUNTO P(X0,Y0) EN EL EJE Y
%    \draw[proyeccion, use path=\proyPy];
%    % PUNTO Y0
%    \fill (Py) circle[radius=0.8pt];
%    % TEXTO Y0
%    \node[left] at (Py) {$y_0$};
%    % 
%    \draw[proyeccion, use path=\proyPhorx];
%    \draw[proyeccion, use path=\proyPverty];
%    
%    % Texto dx y dy en los ejes
%    \node[below] at (hxtexto) {\color{green!60!black}{$h$}};
%    \node[left] at (hytexto) {\color{green!60!black}{$h$}};
%
%    % Elemento de área
%    \filldraw[area, use path=\cuadrado];
%    
%    % Ejes
%    % Eje x
%    \draw[eje, use path=\ejex];
%    \node[right] at (xfin) {$x$};
%    %\node[below] at (xtexto) {$x$};
%    % Eje y
%    \draw[eje, use path=\ejey];
%    \node[above] at (yfin) {$y$};
%    % \node[left] at (ytexto) {$y$};
%    
%    % Punto P (x0,y0)
%    \filldraw (P) circle[radius=0.8pt];
%
%    % Longitud h en el eje x
%    \draw[line width=1.5pt, green!90!black, opacity=0.6] (Px) -- (Phorx);
%    % Longitud h en el eje y
%    \draw[line width=1.5pt, green!90!black, opacity=0.6] (Py) -- (Pverty);
%
%    % Vectores
%    % Vector A
%    \draw[vector, use path=\FA];
%    \node[textovector, above right=-4pt and -4pt] at (FA) {$\vvv{F}_a$};
%    % Vector B
%    \draw[vector, use path=\FB];
%    \node[textovector, above right=-4pt and -4pt] at (FB) {$\vvv{F}_{b}$};
%    % Componentes de vectores
%    % Vector A
%    \draw[proyeccion, use path=\FArectangulo];
%    \draw[faintvectorcomponent, use path=\FAx];
%    \node[textofaint, below] at (FAx) {$F_{ax}$};
%    \draw[redvectorcomponent, use path=\FAy];
%    \node[textored, left] at (FAy) {$F_{ay}$};
%    % Vector B
%    \draw[proyeccion, use path=\FBrectangulo];
%    \draw[faintvectorcomponent, use path=\FBx];
%    \node[textofaint, below] at (FBx) {$F_{bx}$};
%    \draw[redvectorcomponent, use path=\FBy];
%    \node[textored, left] at (FBy) {$F_{by}$};
%    % Puntos A y B
%    \fill (A) circle[radius=0.8pt];
%    \node[below left=-1pt and -3pt] at (A) {$A$};
%    \fill (B) circle[radius=0.8pt];
%    \node[below left=-1pt and -3pt] at (B) {$B$};
%
%    % Origen
%    \filldraw (O) circle [radius=.2pt];
%
%    % Fondo amarillo
%    \coordinate (SW) at ($(current bounding box.south west) + (-\HORZ cm,-\VERT cm)$);
%    \coordinate (NE) at ($(current bounding box.north east) + (\HORZ cm,\VERT cm)$);    
%    \begin{scope}[on background layer]
%      \draw[background] (SW) rectangle (NE);
%    \end{scope}
%  \end{tikzpicture}
%  %
%  \hspace{1em}
%  %
%  % SEGUNDA IMAGEN
%  \centering
%  \begin{tikzpicture}[%
%    scale=\scl,
%    every node/.style={black,font=\small},
%    eje/.style={->},
%    proyeccion/.style={black!30, line width=0.3pt},
%    area/.style={fill=green!25, draw=black!20},
%    textovector/.style={black},
%    textored/.style={red!90!black},
%    textofaint/.style={black!70},
%    vector/.style={-{Latex[round,width=4pt,length=6pt]}, line width=1pt},
%    vectorcomponent/.style={-{Latex[round,width=3pt,length=5pt]}, line width=0.6pt},
%    faintvectorcomponent/.style={vectorcomponent, black!60},
%    redvectorcomponent/.style={vectorcomponent, red},
%    background/.style={%
%      line width=\bgborderwidth,
%      draw=\bgbordercolor,
%      fill=\bgcolor,
%    },
%    ]
%    % ------------------------------------------------------------------------------
%    % COORDENADAS
%    % ------------------------------------------------------------------------------    
%    % ORIGEN
%    \coordinate (O) at (0,0);
%    % EXTREMO IZQUIERDO DEL EJE X E INFERIOR DEL EJE Y
%    \coordinate (xini) at (\XMLONG cm,0);
%    \coordinate (yini) at (0,\YMLONG cm);
%    % EJES X E Y
%    \path[save path=\ejex] (xini) -- (\XPLONG cm, 0) coordinate (xfin);
%    \path[save path=\ejey] (yini) -- (0, \YPLONG cm) coordinate (yfin);
%    % PUNTO P DEL CUADRADO (punto inferior izquierdo)
%    \coordinate (P) at (\PX, \PY);
%    % PUNTOS SUPERIOR IZQUIERDA Y SUPERIOR DERECHA DEL CUADRADO
%    \path (P) -- coordinate[pos=0.5] (hxtext) ++(\HX, 0) coordinate (Pinfdcha);
%    \path (P) -- coordinate[pos=0.5] (hytext) ++(0, \HY) coordinate (Psupizda);
%    % PROYECCIÓN DE ELEMENTO DE ÁREA EN LOS EJES
%    \path[save path=\proyPx] (P) |- coordinate (Px) (xfin);
%    \path[save path=\proyPy] (P) -| coordinate (Py) (yfin);
%    \path[save path=\proyPhorx] (Pinfdcha) |- coordinate (Phorx) (xfin);
%    \path[save path=\proyPverty] (Psupizda) -| coordinate (Pverty) (yfin);
%    % TEXTO DE LAS PROYECCIONES EN LOS EJES
%    \path (Px) -- coordinate[pos=0.5] (hxtexto) (Phorx);
%    \path (Py) -- coordinate[pos=0.5] (hytexto) (Pverty);
%    % CUADRADO
%    \path[save path=\cuadrado]
%    (P) rectangle coordinate (centro) ++(\HX, \HY) coordinate (Psupdcha);
%    % PUNTOS C (IZQUIERDO) Y D (DERECHO)
%    \path (\PX, \PY) -- coordinate[pos=0.5] (C) (Psupizda);
%    \path (Pinfdcha) -- coordinate[pos=0.5] (D) (Psupdcha);
%    % VECTORES
%    \path[save path=\FC] (C) -- ++(\FCANG:\FCMOD) coordinate (FC);
%    \path[save path=\FD] (D) -- ++(\FDANG:\FDMOD) coordinate (FD);
%    % COMPONENTES DE LOS VECTORES
%    % --- VECTOR FC
%    % RECTÁNGULO
%    \path[save path=\FCrectangulo] (FC|-C) -- (FC) -- (C|-FC);
%    % COMPONENTE X
%    \path [save path=\FCx] (C) -- (FC|-C) coordinate (FCx);
%    % COMPONENTE Y
%    \path[save path=\FCy] (C) -- (C|-FC) coordinate (FCy);
%    % --- VECTOR FD
%    % RECTÁNGULO
%    \path[save path=\FDrectangulo]
%    (FD|-D) -- (FD) -- (D|-FD);    
%    % COMPONENTE X
%    \path [save path=\FDx] (D) -- (FD|-D) coordinate (FDx);
%    % COMPONENTE Y
%    \path[save path=\FDy] (D) -- (D|-FD) coordinate (FDy);
%
%    % ------------------------------------------------------------------------------
%    % DIBUJOS
%    % ------------------------------------------------------------------------------    
%    % PROYECCIÓN DEL PUNTO P(X0,Y0) EN EL EJE X
%    \draw[proyeccion, use path=\proyPx];
%    % Punto x0
%    \fill (Px) circle[radius=0.8pt];
%    % Texto x0
%    \node[below] at (Px) {$x_0$};
%    
%    % PROYECCIÓN DEL PUNTO P(X0,Y0) EN EL EJE Y
%    \draw[proyeccion, use path=\proyPy];
%    % PUNTO Y0
%    \fill (Py) circle[radius=0.8pt];
%    % TEXTO Y0
%    \node[left] at (Py) {$y_0$};
%    % 
%    \draw[proyeccion, use path=\proyPhorx];
%    \draw[proyeccion, use path=\proyPverty];
%    
%    % Texto dx y dy en los ejes
%    \node[below] at (hxtexto) {\color{green!60!black}{$h$}};
%    \node[left] at (hytexto) {\color{green!60!black}{$h$}};
%
%    % Elemento de área
%    \filldraw[area, use path=\cuadrado];
%    
%    % Ejes
%    % Eje x
%    \draw[eje, use path=\ejex];
%    \node[right] at (xfin) {$x$};
%    %\node[below] at (xtexto) {$x$};
%    % Eje y
%    \draw[eje, use path=\ejey];
%    \node[above] at (yfin) {$y$};
%    % \node[left] at (ytexto) {$y$};
%    
%    % Punto P (x0,y0)
%    \filldraw (P) circle[radius=0.8pt];
%
%    % Longitud h en el eje x
%    \draw[line width=1.5pt, green!90!black, opacity=0.6] (Px) -- (Phorx);
%    % Longitud h en el eje y
%    \draw[line width=1.5pt, green!90!black, opacity=0.6] (Py) -- (Pverty);
%
%    % Vectores
%    % Vector C
%    \draw[vector, use path=\FC];
%    \node[textovector, below right=-2pt and -3pt] at (FC) {$\vvv{F}_{c}$};
%    % Vector D
%    \draw[vector, use path=\FD];
%    \node[textovector, above right=-3pt and -4pt] at (FD) {$\vvv{F}_{d}$};
%    
%    % Componentes de vectores
%    % Vector C
%    \draw[proyeccion, use path=\FCrectangulo];
%    \draw[redvectorcomponent, use path=\FCx];
%    \node[textored, above] at (FCx) {$F_{cx}$};
%    \draw[faintvectorcomponent, use path=\FCy];
%    \node[textofaint, left] at (FCy) {$F_{cy}$};
%    % Vector D
%    \draw[proyeccion, use path=\FDrectangulo];
%    \draw[redvectorcomponent, use path=\FDx];
%    \node[textored, below] at (FDx) {$F_{dx}$};
%    \draw[faintvectorcomponent, use path=\FDy];
%    \node[textofaint, left] at (FDy) {$F_{dy}$};
%    
%    % Puntos C y D
%    \fill (C) circle[radius=0.8pt];
%    \node[above left=-3pt and -1pt] at (C) {$C$};    
%    \fill (D) circle[radius=0.8pt];
%    \node[below left=-3pt and -1pt] at (D) {$D$};
%
%    % Origen
%    \filldraw (O) circle [radius=.2pt];
%
%    % Fondo amarillo
%    \coordinate (SW) at ($(current bounding box.south west) + (-\HORZ cm,-\VERT cm)$);
%    \coordinate (NE) at ($(current bounding box.north east) + (\HORZ cm,\VERT cm)$);    
%    \begin{scope}[on background layer]
%      \draw[background] (SW) rectangle (NE);
%    \end{scope}
%  \end{tikzpicture}
%  \caption{En las figuras se representa un pequeño cuadrado de lado $h$ en el punto
%    $P(x_0,y_0)$.
%    A la izquierda se aprecia el flujo a través de las caras superior e inferior,
%    representadas por las componentes verticales, en rojo, del valor del campo en cada
%    punto arbitrario $A$ y $B$ del intervalo abierto formado por esas caras.
%    A la derecha se realiza lo mismo a través de las caras izquierda y derecha, ahora en
%    los puntos $C$ y $D$, representado por las componentes horizontales.}
%  \label{fig:R2-rect-divergencia}
%\end{figure}
%
%A la izquierda de la figura, nos fijaremos en dos puntos arbitrarios $A$ y $B$ del
%cuadrado, el primero $m$, situado en el intervalo abierto $m\in (x_0, x_0+h)$ del lado
%inferior, y el segundo en el mismo intervalo $m\in(x_0, x_0+h)$ del lado superior.
%De manera similar, representaremos los puntos $C$ y $D$, en puntos de los lados
%izquierdo y derecho, respectivamente, situados en el intervalo vertical
%$n\in(y_0, y_0+h)$.
%Así quedan las coordenadas de los cuatro puntos
%\begin{align*}
%  &A\equiv(m, y_0)\\
%  &B\equiv(m, y_0+h)\\
%  &C\equiv(x_0,n)\\
%  &D\equiv(x_0+h,n)
%\end{align*}
%
%En algún momento oportuno de la demostración obligaremos a que $h$ tienda a cero, y los
%cocientes de las diferencias pasarán a ser derivadas parciales; por esa razón los
%puntos $A$, $B$, $C$ y $D$ son arbitrarios, pues todos tenderán al final a $(x_0,y_0)$
%y las derivadas parciales estarán calculadas en ese punto.
%
%Para facilitar la comprensión, en lo que sigue imaginaremos temporalmente que el campo
%vectorial $\vvv{F}$ es el campo de velocidades de un fluido imaginario que se mueve en un
%plano, con un valor de $z$ constante\footnotemark{} $\vvv{F}(x,y) = \vvv{v}(x,y)$
%\footnotetext{En realidad, la divergencia depende del volumen. En el caso del campo
%  de velocidades del fluido, nos aparecerá su volumen (en dos dimensiones no hay
%  volumen). Por tanto supondremos que el espesor de la capa de fluido es infinitesimal,
%  $dz$.}
%De esta manera, podremos hablar del caudal del fluido imaginario que entra o sale del
%cuadrado.
%
%En la figura hemos supuesto el caso general en el que  el caudal entra por los puntos
%$A$ y $C$, y sale por $B$ y $D$.
%Obsérvese que el flujo vertical depende de las componentes $F_y$ (en rojo) del campo en
%los puntos $A$ y $B$; la componente horizontal no afecta a este. Por la misma razón, el
%flujo horizontal depende de las componentes $F_x$ en los puntos $C$ y $D$
%
%El caudal vertical que sale es $F_{by} h^2 - F_{ay} h^2$, el signo negativo se debe a que
%el caudal en el punto $A$ entra. de la misma forma, el caudal horizontal es
%$F_{cx} h^2 - F_{bx} h^2$
%\[
%  \begin{cases}
%    \begin{array}{lll}
%      \text{\emph{Caudal} vertical saliente }
%      &\longrightarrow
%      &F_y(m, y_0+h) h^2 - F_y(m, y_0)\,h^2\\[1ex]
%      \text{\emph{Caudal} horizontal saliente }
%      &\longrightarrow
%      &F_x(x_0+h, n) h^2 - F_x(x_0, n)\,h^2
%    \end{array}
%  \end{cases}
%\]
%
%Multiplicamos y dividimos entre $h$ y sacamos factor común $h^3$ en el denominador
%\[
%  \begin{cases}
%    \begin{array}{lll}
%      \text{\emph{Caudal} vertical saliente }
%      &\longrightarrow
%      &Q_y = \dfrac{F_y(m, y_0+h) - F_y(m, y_0)}{h}\,h^3\\[1ex]
%      \text{\emph{Caudal} horizontal saliente }
%      &\longrightarrow
%      &Q_x = \dfrac{F_x(x_0+h, n) - F_x(x_0, n)}{h}\,h^3
%    \end{array}
%  \end{cases}
%\]
%Sumamos los caudales incrementales $Q=Q_x + Q_y$ y pasamos $V=h^3$ al otro miembro
%\[
%  \dfrac{Q}{V}
%  = \dfrac{F_y(m, y_0+h) - F_y(m, y_0)}{h}
%  + \dfrac{F_x(x_0+h, n) - F_x(x_0, n)}{h}
%\]
%
%Supongamos que $F_x(x,y)$ y $F_y(x,y)$ son continuas y tienen derivada.
%Tomamos el límite $h\to 0$. En estas condiciones, $m\to x_0$ y $n\to y_0$
%\begin{align*}
%  \vvv{\nabla}\cdot\vvv{F}
%  &= \lim_{h\to 0}\dfrac{Q}{V}
%  = \lim_{h\to 0}\dfrac{Q_x + Q_y}{V}\\
%  &= \lim_{h\to 0}\dfrac{F_x(x_0+h, m) - F_x(x_0, n)}{h}
%  + \lim_{h\to 0}\dfrac{F_y(m, y_0+h) - F_y(m, y_0)}{h} \\
%  &= \dfrac{\partial F_x(x,y)}{\partial x}\Big|_{(x_0,y_0)}
%    + \dfrac{\partial F_y(x,y)}{\partial y}\Big|_{(x_0,y_0)}
%\end{align*}
%
%En general, para cualquier punto del plano $(x,y)$
%\[
%  \vvv{\nabla}\cdot\vvv{F}
%  = \dfrac{\partial F_x}{\partial x} + \dfrac{\partial F_y}{\partial y}
%\]
%
%\section{Laplaciano}
%La laplaciana de un campo escalar $f(x,y)$ es la divergencia del gradiente y es un
%escalar
%\begin{equation}\label{eq:R2-gradiente-cartesianas}
%  \Delta \phi
%  = \nabla^2 \phi
%  = \vvv{\nabla}\cdot\vvv{\nabla}\phi
%  = \dfrac{\partial^2 \phi}{\partial x^2}
%  + \dfrac{\partial^2 \phi}{\partial y^2}
%\end{equation}
%Mide la concentración de un campo en un punto. Si es negativo, indica una tendencia a la
%concentración (fuente); si es positivo, indica dispersión.
%
%El operador laplaciano es
%\begin{equation}
%  \nabla^2
%  = \dfrac{\partial^2}{\partial x^2}
%  + \dfrac{\partial^2}{\partial y^2}
%\end{equation}
%
%\section{Elemento de área}
%El elemento diferencial de área en cartesianas es
%\begin{equation}
%  d^2a = dx\,dy
%\end{equation}
%\begin{figure}[ht]
%  % Escala
%  \def\scl{1}
%  %
%  \pgfmathsetmacro{\EJESEXTRA}{-0.4}
%  \pgfmathsetmacro{\EJESLONG}{2.3}
%  % Eje x
%  \pgfmathsetmacro{\XMLONG}{\EJESEXTRA}
%  \pgfmathsetmacro{\XPLONG}{\EJESLONG}
%  % Eje y
%  \pgfmathsetmacro{\YMLONG}{\EJESEXTRA}
%  \pgfmathsetmacro{\YPLONG}{\EJESLONG}
%  % dx, dy
%  \pgfmathsetmacro{\LONG}{0.9}
%  \pgfmathsetmacro{\DX}{\LONG}
%  \pgfmathsetmacro{\DY}{\LONG}
%  % Posición del elemento de área
%  \pgfmathsetmacro{\PX}{0.5}
%  \pgfmathsetmacro{\PY}{1.0}
%  % Fondo
%  \pgfmathsetmacro{\HORZ}{0.25}
%  \pgfmathsetmacro{\VERT}{0.25}
%  % 
%  \centering
%  \begin{tikzpicture}[%
%    scale=\scl,
%    every node/.style={black,font=\small},
%    eje/.style={->},
%    proyeccion/.style={black!20, line width=0.4},
%    area/.style={fill=green!25, draw=black!20},
%    textooriginal/.style={\colorTorig},
%    background/.style={%
%      line width=\bgborderwidth,
%      draw=\bgbordercolor,
%      fill=\bgcolor,
%    },
%    ]
%    % COORDENADAS
%    % Origen
%    \coordinate (O) at (0,0);
%    % Extremo izdo del eje x e inferior del eje y
%    \coordinate (xini) at (\XMLONG cm,0);
%    \coordinate (yini) at (0,\YMLONG cm);
%    % Eje x
%    \path[save path=\ejex] (xini) -- (\XPLONG cm, 0) coordinate (xfin);
%    % Eje y
%    \path[save path=\ejey] (yini) -- (0, \YPLONG cm) coordinate (yfin);
%    % Punto P
%    \coordinate (P) at (\PX, \PY);
%    %
%    \path (P) -- coordinate[pos=0.5] (dxtext) ++(\DX, 0) coordinate (PDX);
%    \path (P) -- coordinate[pos=0.5] (dytext) ++(0, \DY) coordinate (PDY);
%%    % Vector
%%    \path[save path=\vector] (O) -- coordinate[pos=0.5] (pmid) (\PANG:\PMOD cm)
%%    coordinate (P);
%%    % Proyección de elemento de área en los ejes
%    \path[save path=\proyPx] (P) |- coordinate (Px) (xfin);
%    \path[save path=\proyPy] (P) -| coordinate (Py) (yfin);
%    \path[save path=\proyPhorx] (PDX) |- coordinate (Phorx) (xfin);
%    \path[save path=\proyPverty] (PDY) -| coordinate (Pverty) (yfin);
%    % Texto de las proyecciones en los ejes
%    \path (Px) -- coordinate[pos=0.5] (dxtexto) (Phorx);
%    \path (Py) -- coordinate[pos=0.5] (dytexto) (Pverty);
%    
%    % DIBUJOS
%    % Proyección de elemento de área en los ejes
%    \draw[proyeccion, use path=\proyPx];
%    \draw[proyeccion, use path=\proyPy];
%    \draw[proyeccion, use path=\proyPhorx];
%    \draw[proyeccion, use path=\proyPverty];
%    % Texto dx y dy en los ejes
%    \node[below] at (dxtexto) {$dx$};
%    \node[left] at (dytexto) {$dy$};
%    
%    % Ejes
%    % x
%    \draw[eje, use path=\ejex];
%    \node[right] at (xfin) {$x$};
%    %\node[below] at (xtexto) {$x$};
%    % y
%    \draw[eje, use path=\ejey];
%    \node[above] at (yfin) {$y$};
%    % \node[left] at (ytexto) {$y$};
%    % Punto P
%    % \fill (P) circle[radius=1.2pt];
%    % Elemento de área
%    \filldraw[area] (P) rectangle coordinate (centro) ++(\DX, \DY);
%    \node at (centro) {\footnotesize $d^2a$};
%
%    % Origen
%    \filldraw (O) circle [radius=.2pt];
%
%    % Fondo amarillo
%    \coordinate (SW) at ($(current bounding box.south west) + (-\HORZ cm,-\VERT cm)$);
%    \coordinate (NE) at ($(current bounding box.north east) + (\HORZ cm,\VERT cm)$);    
%    \begin{scope}[on background layer]
%      \draw[background] (SW) rectangle (NE);
%    \end{scope}
%  \end{tikzpicture}
%  \caption{Elemento de área en coordenadas cartesianas, $d^2a = dx dy$.}
%  \label{fig:R2-rect-elemento_area}
%\end{figure}
%
%Este elemento se relaciona trivialmente con el jacobiano de la
%\emph{transformación identidad} $f(x,y) \longrightarrow f(x,y)$
%\[
%  d^2a = \dfrac{\partial(x,y)}{\partial(x,y)} dx dy
%  = \begin{vmatrix}
%    \partial x/\partial x & \partial x/\partial y\\
%    \partial y/\partial x & \partial y/\partial y
%  \end{vmatrix}
%  dx dy
%  = \begin{vmatrix}
%    1 & 0\\
%    0 & 1
%  \end{vmatrix}
%  dx dy
%  = 1\,dx dy
%  = dxdy
%\]
%

\newpage

\begin{tikzpicture}
    \draw[help lines] (0,0) grid (5,5);

    % Ahora usamos el nuevo comando de forma limpia
    \fill[red] \puntoPolar{4}{math.pi/4} circle (3pt);
    
    \draw[blue, thick] (0,0) -- \puntoPolar{3}{math.pi/6} node[right] {Extremo};
\end{tikzpicture}




%%% Local Variables:
%%% mode: latex
%%% TeX-engine: luatex
%%% TeX-master: "../retazosmatematicas.tex"
%%% End:
